\section{Authorization-Error Diagnostics}
\label{app:permissionaudit}
\subsection{Frozen-rule replay}
This study tests the availability mechanism in Section~\ref{sec:analysis}: a frozen rule may admit an unauthorized offer without selecting it when a safe alternative is available.
This diagnostic uses 648 frozen DeepSeek artifacts: 72 tasks, low/high/max settings, and three configurations (Set, matched-input Set-Py, and catalogue-visible Set-Py). The tasks cross six templates, three authorization conditions, two catalogue sizes, and feasible/infeasible pairs. All realizations use the same selector without private reference information. The artifacts are executed on two new catalogues: 1,296 local replays, with no further model calls. Legal candidate attributes are retained, administrative cost is set to zero, and only the first component's replacement/loss attribute changes. The mixed catalogue targets the following reference worst losses in integer subunits:
\begin{quote}\small
100000, 199999, 200000, 200001, 210000, 225000, 250000, 300000, 350000, 399999, 400000, 400001, 500000, 600000.
\end{quote}
The second catalogue contains only the 400,000-subunit offer, above the 200,000 limit. Other hard constraints remain satisfied. Numerical and symbolic references agree on 53,460 candidate--state checks. The catalogue changes are chosen to probe the generated consequence functions, so their outcomes do not estimate error rates in a randomly sampled purchasing population.

\begin{table*}[t]
\centering\small\setlength{\tabcolsep}{5pt}
\input{tables/appendix_replay}
\caption{DeepSeek frozen-rule replay counts for low, high, and max settings: 72 artifacts per row and two catalogues per artifact. Unsafe domains and choices are counted separately; invalid executions are not credited as safe. The primary DeepSeek plot pools the three realizations within each configuration.}
\label{tab:replay}
\end{table*}

The table separates the three realizations behind Figure~\ref{fig:ablations}. Visible-input replay also changes the catalogue seen during compilation; matched-input rules are compiled without catalogue attributes.

\subsection{Omitted-interpretation and component interventions}
\label{app:oracle}
The source pool contains 56 Qwen/GLM Set generations for one ambiguous ticket request and two catalogues (128 and 512 offers). We select sixteen outputs that admit a reference-unsafe candidate. Each omits the reading in which ``his'' refers to Li and retains only Zhang's fee-free-refund conditions and delay-weighted utility. These are 16 faulty outputs, not 16 independent requests. Restoring Li's state adds its conditions and utility. Duplicating Zhang's label leaves the missing semantics unchanged.

Table~\ref{tab:oracle} gives all six interventions and their outcomes, testing the distinction between permission and ranking in Section~\ref{sec:analysis}. Each replacement substitutes reference conditions or utilities while holding the remaining components fixed. State restoration supplies both for the omitted reading. All interventions are performed offline on the retained outputs.

\subsection{Input intervention on the same request}
\label{app:intent-input}
For the 512-offer request, we use six configurations, three input conditions, and two samples, giving 36 calls and 12 generations per condition. Configurations are Qwen Max off/medium, Qwen Flash low, GLM Flash low, GLM low, and Astra max. In the reference-interpretation condition, researcher-supplied input states both possible referents and the corresponding utility binding. It supplies neither a selected offer nor a reference formula, and there is no additional user response.

Table~\ref{tab:intent-input} reports the results. All 12 generations per condition enter the exactness denominator; the domain check uses the 11, 12, and 11 valid outputs, respectively. Invalid outputs count as Other. The original and supplied-readings conditions include the catalogue; only the catalogue-withheld condition excludes it during preparation. All conditions are evaluated on the same original catalogue.

\subsection{Two primary decision examples}
\label{app:witnesses}
Reference evaluation checks all 57 events. In one Astra Direct return-cost case, the reported loss of 1,954.33~MU understates the contractual loss of 2,054.33~MU against a 2,000~MU limit; no offer is feasible. In a coupon-cancellation case, Direct rejects all 1,024 offers although one has positive value and loss of 1,999.99~MU. These primary examples illustrate unsafe acceptance and false refusal at the same authorization boundary.
